% CSC 898 Master's Thesis — First Draft
% San Francisco State University, Department of Computer Science
% Single-file draft. Uses 1-inch margins per Graduate Studies guidance.
% No official department LaTeX template was found in the repository.
% Fill administrative placeholders before format check.

\documentclass[12pt,oneside]{report}

\usepackage[margin=1in]{geometry}
\usepackage[T1]{fontenc}
\usepackage[utf8]{inputenc}
\usepackage{microtype}
\usepackage{setspace}
\usepackage{booktabs}
\usepackage{amsmath}
\usepackage{graphicx}
\usepackage{hyperref}

\hypersetup{colorlinks=true,linkcolor=black,citecolor=black,urlcolor=black}

\newcommand{\thesisauthor}{[AUTHOR NAME]}
\newcommand{\thesisadvisor}{[ADVISOR NAME]}
\newcommand{\thesisdegree}{Master of Science in Computer Science}
\newcommand{\thesisuniversity}{San Francisco State University}
\newcommand{\thesissubmissiondate}{[SUBMISSION SEMESTER AND YEAR]}

\title{SkillLLM\_v2: Context-Aware Graph-Augmented Skill Extraction\\and Career Intelligence from Employment Text}
\author{\thesisauthor}
\date{\thesissubmissiondate}

\begin{document}
\onehalfspacing

% ---- Title page ----
\begin{titlepage}
  \centering
  {\large \thesisuniversity\par}
  {\large Department of Computer Science\par}
  \vspace{2cm}
  {\LARGE\bfseries SkillLLM\_v2: Context-Aware Graph-Augmented Skill Extraction and Career Intelligence from Employment Text\par}
  \vspace{2cm}
  {\large A Thesis Submitted in Partial Fulfillment of the Requirements for the Degree\par}
  {\large \thesisdegree\par}
  \vspace{2cm}
  {\large by\par}
  {\large \thesisauthor\par}
  \vspace{1.5cm}
  {\large San Francisco, California\par}
  {\large \thesissubmissiondate\par}
  \vfill
  {\large Committee Chair: \thesisadvisor\par}
\end{titlepage}

% ---- Abstract ----
\chapter*{Abstract}
\addcontentsline{toc}{chapter}{Abstract}

Employment-related text encodes skills as unstructured spans rather than normalized database entries.
This thesis presents SkillLLM\_v2, a locally executed skill intelligence system built on supervised generative extraction from open-weight language models.
The work proceeds in two verified stages documented in the SkillLLM repository.

First, a controlled SkillSpan benchmark compared five instruction-tuned models under identical splits, prompts, LoRA/4-bit training, and evaluation code.
Mistral Small 3.1 24B achieved the strongest completed result among feasible models (exact span F1 0.665 on 3{,}569 test examples).
Gemma 4 26B failed repeatedly on available 2$\times$24\,GB GPU hardware and is reported as hardware-infeasible rather than ranked.

Second, a Stage-2 system centers on an 80-epoch Mistral LoRA extractor (exact span F1 0.670), augmented by GLiNER and dictionary candidates, boundary refinement, taxonomy linking, and a non-regression wrapper gate before graph-based job matching and course recommendation.
Encoder baselines under the same evaluation protocol reached exact span F1 of 0.059 (BERT) and 0.076 (zero-shot GLiNER).
The full context-aware graph wrapper preserved anchor extraction quality while achieving taxonomy linking coverage of 0.890.
Labeled evaluation of job matching and course recommendation was not available in repository artifacts; those modules are demonstrated on synthetic demo inputs only.

\tableofcontents
\listoftables

% ============================================================
\chapter{Introduction}
\label{ch:intro}

\section{Problem and Motivation}

Resumes, job postings, and professional profiles contain explicit mentions of technical and interpersonal skills.
Downstream applications---job matching, labor-market analytics, and training recommendations---require structured skill representations with traceable evidence in the source text~\cite{senger2024survey,otani2025hrsurvey}.
Skill extraction from English job postings has been formalized as a span labeling task distinguishing hard and soft skills~\cite{zhang2022skillspan}.
Recent work shows that repurposing general-purpose large language models through supervised adaptation can improve extraction quality relative to classical sequence taggers and extreme multi-label skill classifiers~\cite{herandi2024skillllm,nguyen2024rethinking,zhang2024nnose}.

This thesis documents what was actually implemented, trained, and measured in the SkillLLM research repository.
It does not describe planned components as completed work.

\section{Research Questions}

The completed experiments support the following questions:

\begin{enumerate}
  \item How do contemporary open-weight LLMs compare when fine-tuned on SkillSpan under a shared protocol?
  \item Which backbone is strongest among models feasible on 2$\times$ NVIDIA RTX 3090 (24\,GB) GPUs?
  \item How do encoder baselines compare with a fine-tuned generative Mistral extractor?
  \item Can a guarded hybrid wrapper improve or preserve span metrics while enabling taxonomy normalization and downstream reasoning?
  \item What error modes remain after fine-tuning and wrapper processing?
\end{enumerate}

Downstream recommendation performance (RQ4 extension) is not quantitatively answered because labeled job/course evaluation sets were absent from verified artifacts.

\section{Contributions}

\begin{enumerate}
  \item A reproducible multi-model SkillSpan benchmark pipeline comparing five open-weight LLMs with shared preprocessing, QLoRA training, JSON decoding, span alignment, and multi-metric evaluation.
  \item Empirical selection of Mistral Small 3.1 24B as the strongest completed benchmark model, with documented Gemma 4 26B failure on available hardware.
  \item SkillLLM\_v2 Stage-2: a context-aware graph-augmented architecture with hybrid candidate extraction, non-regression gating, taxonomy linking, and NetworkX-based relational reasoning.
  \item Verified comparative metrics on SkillSpan test data for BERT, GLiNER, raw Mistral LoRA, and wrapper variants.
\end{enumerate}

\section{Thesis Organization}

Chapter~\ref{ch:background} reviews related work.
Chapter~\ref{ch:method} defines the task and metrics.
Chapter~\ref{ch:data} describes datasets.
Chapter~\ref{ch:arch} presents the architecture.
Chapters~\ref{ch:setup}--\ref{ch:results} cover experiments and results.
Chapters~\ref{ch:discuss}--\ref{ch:concl} analyze findings, limitations, and conclusions.

\section{Formatting Note}

This document is prepared for CSC 898 (Master's Thesis) at San Francisco State University.
Graduate Studies requires department-approved formatting, 1-inch margins, and accessibility review before library submission.
Final approval by the Graduate Program Specialist is required.

% ============================================================
\chapter{Background and Related Work}
\label{ch:background}

\section{Skill Extraction and Labor-Market NLP}

SkillSpan provides expert-annotated hard and soft skill spans on English job posting sentences (4{,}800 train, 3{,}174 dev, 3{,}569 test examples) under CC BY 4.0 licensing~\cite{zhang2022skillspan}.
It remains the primary evaluation benchmark in this thesis.
Senger et al.\ survey deep learning approaches to skill extraction and classification from job postings and identify these tasks as central to computational job-market analysis~\cite{senger2024survey}.

Earlier systems addressed skill identification and normalization at scale~\cite{zhao2015skill,hoang2018largescale} and implicit skill recovery for recommendation~\cite{gugnani2020implicit}.
Extreme multi-label formulations treat skill retrieval as classification over large label spaces~\cite{bhola2020retrieving,goyal2023jobxmlc}.
ESCOXLM-R and related work incorporate taxonomy-driven pre-training for the job market domain~\cite{zhang2023escoxlmr,levrang2014esco}.
Entity linking and zero-shot ESCO matching extend extraction toward normalized skill inventories~\cite{zhang2024entitylinking,clavie2023batteries}.
Soft-skill requirements pose distinct lexical and matching challenges~\cite{sayfullina2018softskill,calanca2019responsible}.

Recent LLM-centric lines include Skill-LLM~\cite{herandi2024skillllm}, SkillGPT~\cite{li2023skillgpt}, and re-evaluations of LLM skill extraction in the job domain~\cite{nguyen2024rethinking}.
NNOSE proposes nearest-neighbor occupational skill extraction~\cite{zhang2024nnose}, which the repository's earlier resume benchmark also evaluated.
Synthetic job posting generation (JobSkape) supports data augmentation for skill matching research~\cite{magron2024jobskape}.
Multilingual and large-scale skill mapping have moved rapidly at the frontier: SkiLLens maps novel skills from millions of European job ads~\cite{desanto2026skillens}, and recent work extracts skills from resumes and offers across six languages~\cite{vasquezrodriguez2026sixlanguages,kavas2025multilingual}.

\section{Named Entity Recognition and Span Extraction}

Classical sequence labeling with BiLSTM-CRF and CNN-CRF architectures established strong NER baselines~\cite{lample2016ner,ma2016bilstmcnncrf,huang2015bilstmcrf}.
Contextual embeddings from Flair and pooled variants improved sequence taggers~\cite{akbik2018flair,akbik2019pooled}.
Span-oriented pre-training and prediction (SpanBERT, SpanNER) treat entities as span selections rather than token tags~\cite{joshi2020spanbert,fu2021spanner}.
GLiNER generalizes open-label NER with a compact bidirectional transformer~\cite{zaratiana2023gliner}; SkillLLM\_v2 uses it as an auxiliary candidate source.
GPT-NER and related work frame NER through generative LLMs~\cite{wang2023gptner}, aligning with the generative JSON extraction approach used for Mistral in this thesis.
DyGIE++ jointly models entities and relations via contextualized spans~\cite{wadden2019dygiepp}.
Domain-adaptive pre-training remains important when transferring general language models to specialized corpora~\cite{gururangan2020dontstop,gnehm2022transfer}.

\section{Transformers, Matching, and HR NLP}

Transformer encoders~\cite{vaswani2017attention,devlin2019bert,liu2019roberta,conneau2020xlmr} underpin both extraction baselines and matching systems.
Sentence-BERT, SimCSE, and supervised contrastive learning produce semantic embeddings for retrieval and matching~\cite{reimers2019sentencebert,gao2021simcse,khosla2020supervisedcontrastive}.
Person--job fit has been studied with domain-adapted deep matching networks and context-aware transformers~\cite{bian2019domain,li2020competence}.
ConFit v2 improves resume--job matching with hypothetical resume embeddings and hard-negative mining~\cite{yu2025confitv2}.
Human versus LLM resume matching and bias evaluation are active 2025 research areas~\cite{vaishampayan2025humanllm,iso2025bias}.
Graph-based job recommendation over skill embeddings appears in Skills2Job~\cite{giabelli2021skills2job}.
Otani et al.\ provide a broader NLP-for-HR survey~\cite{otani2025hrsurvey}.

\section{Fine-Tuning, Retrieval, and This Thesis}

Parameter-efficient fine-tuning via LoRA and QLoRA enables local adapter training of large models~\cite{hu2022lora,dettmers2023qlora}.
Unified tooling such as LlamaFactory supports reproducible fine-tuning workflows~\cite{zheng2024llamafactory}.
The Mistral and Mixtral model families provide open-weight backbones for instruction-tuned extraction~\cite{jiang2023mistral,jiang2024mixtral}.
Retrieval-augmented generation, dense passage retrieval, and late-interaction retrievers (ColBERT) ground dynamic job and course context~\cite{lewis2020rag,karpukhin2020dpr,khattab2020colbert}.

SkillLLM\_v2 situates itself after this literature chronologically: it selects a fine-tuned Mistral extractor on SkillSpan, adds guarded hybrid candidates and taxonomy linking inspired by ESCO normalization work, and implements graph-structured downstream matching consistent with recent HR matching research---while reporting only metrics verified in the repository.

% ============================================================
\chapter{Methodology}
\label{ch:method}

\section{Task Definition}

Given input text $x$, the system predicts a set of spans $\hat{S} = \{(\hat{s}_j, \hat{\ell}_j)\}$ where each surface form $\hat{s}_j$ is a substring of $x$ and $\hat{\ell}_j \in \{\textsc{Hard\_Skill}, \textsc{Soft\_Skill}\}$.
Gold annotations follow the same schema on SkillSpan~\cite{zhang2022skillspan}.

\section{Generative Extraction}

The primary extractor is Mistral Small 3.1 24B Instruct, an open-weight instruction-tuned model in the Mistral lineage~\cite{jiang2023mistral}, fine-tuned with LoRA adapters~\cite{hu2022lora,dettmers2023qlora} on prompt--JSON target pairs.
Training minimizes causal language modeling loss on the target tokens (supervised fine-tuning in \texttt{train\_sft\_v2.py}), following the generative extraction paradigm explored in Skill-LLM and GPT-NER~\cite{herandi2024skillllm,wang2023gptner}.
This is structured generative extraction, not classical BiLSTM-CRF tagging~\cite{huang2015bilstmcrf} or token-level fusion.

Inference produces JSON skill lists; predicted strings are aligned to $x$ by exact and case-insensitive substring search.
Unaligned predictions are counted separately.

\section{Evaluation Metrics}

Reported metrics from repository evaluation scripts include:
\begin{itemize}
  \item Exact span F1 (boundary and label must match)
  \item Partial span F1 (overlap with correct label)
  \item Token-level F1
  \item Per-label hard/soft skill F1
  \item Parse failure rate and unaligned prediction rate
  \item Taxonomy linking coverage (Stage-2 wrapper)
\end{itemize}

\section{Two Experimental Phases}

\textbf{Phase A (legacy benchmark):} five LLMs, three epochs each, metrics in \texttt{SkillLLM\_v2/results/skillspan\_baselines/}.

\textbf{Phase B (Stage-2):} Mistral-focused 80-epoch training, encoder baselines, wrapper ablations, metrics in \texttt{SkillLLM\_v2\_stage2/artifacts/skillllm\_v2/}.

Results from these phases are labeled explicitly when compared; they used different epoch counts and LoRA ranks.

% ============================================================
\chapter{Datasets and Preprocessing}
\label{ch:data}

\section{SkillSpan}

The primary benchmark uses SkillSpan exported to \texttt{data/source/skillspan\_hf/}.
The benchmark pipeline converts records to JSONL with \texttt{text} and \texttt{skills} fields (character offsets included in gold data).
Table~\ref{tab:splits} lists verified split sizes from \texttt{summary\_report.json}.

\begin{table}[htbp]
  \centering
  \caption{SkillSpan split sizes.}
  \label{tab:splits}
  \begin{tabular}{lr}
    \toprule
    Split & Examples \\
    \midrule
    Train & 4{,}800 \\
    Dev   & 3{,}174 \\
    Test  & 3{,}569 \\
    \bottomrule
  \end{tabular}
\end{table}

Generative training uses a fixed system prompt and user prompt requiring JSON output with \textsc{Hard\_Skill} or \textsc{Soft\_Skill} labels; training targets omit offset fields.
SkillSpan converted examples follow the same span-and-label interface used in prior SkillSpan evaluations~\cite{zhang2022skillspan} and related job-market corpora~\cite{green2022benchmark,gnehm2022finegrained}.

\section{Taxonomies and Normalization Targets}

Stage-2 linking targets ESCO-style alias resources~\cite{levrang2014esco,zhang2024entitylinking} via custom alias files in the repository.
Full ESCO/O*NET ingestion is an extension point; reported linking coverage (0.890) used the configured custom taxonomy.

\section{Other Repository Data}

\texttt{Extract\_jobDesc/} contains a separate BERT NER track~\cite{devlin2019bert} on scraped job descriptions (entity F1 0.462 gold-only; 0.710 silver-train/gold-eval on verified \texttt{test\_metrics.json} files).
\texttt{Experiments/skill\_extraction\_benchmark/} contains an earlier resume-sentence benchmark with different metrics (e.g., NNOSE span F1 0.535 per \texttt{comparison.md})~\cite{zhang2024nnose}; it is not merged into SkillSpan tables because metric definitions differ.

Stage-2 demo inputs for job matching use \texttt{inputs/demo/jobs.jsonl} without gold relevance labels.

% ============================================================
\chapter{SkillLLM\_v2 Architecture}
\label{ch:arch}

Stage-2 implements a modular pipeline (\texttt{src/skillllm\_v2/}):

\begin{enumerate}
  \item \textbf{Mistral LoRA extractor} --- primary generative span extractor (\texttt{mistral\_extractor.py}), following the LLM extraction line~\cite{herandi2024skillllm,nguyen2024rethinking}.
  \item \textbf{GLiNER extractor} --- lightweight recall booster~\cite{zaratiana2023gliner} (zero-shot in verified runs).
  \item \textbf{Dictionary extractor} --- exact/fuzzy alias matching against taxonomy files~\cite{zhao2015skill,clavie2023batteries}.
  \item \textbf{Candidate merger} --- deduplication with provenance preservation.
  \item \textbf{Boundary refiner} --- rule-based span cleanup.
  \item \textbf{Span verifier and confidence calibrator} --- rule-based with extension points for trained verifiers.
  \item \textbf{Taxonomy linker} --- maps spans to canonical skill IDs~\cite{zhang2024entitylinking,levrang2014esco}.
  \item \textbf{Graph layer} --- NetworkX nodes/edges for users, jobs, courses, skills (cf.\ graph job recommenders~\cite{giabelli2021skills2job,goyal2023jobxmlc}).
  \item \textbf{Job matcher, gap detector, course recommender} --- scoring modules with explainable outputs, related to person--job fit literature~\cite{bian2019domain,yu2025confitv2,li2020competence}.
\end{enumerate}

Each skill record carries \texttt{surface\_form}, \texttt{label}, offsets, \texttt{source}, \texttt{confidence}, and \texttt{evidence}.

The \textbf{non-regression gate} compares wrapper metrics to the Raw Mistral anchor; if guardrails fail, extraction metrics fall back to anchor predictions while linking metadata may still be recorded.

Stable profile and taxonomy summaries are cached in context packets; dynamic job text is retrieved at inference time~\cite{lewis2020rag,karpukhin2020dpr,reimers2019sentencebert}.
Graph retrieval explains paths from user skills to missing skills and courses.

% ============================================================
\chapter{Experimental Setup}
\label{ch:setup}

\section{Hardware and Software}

Verified Stage-2 environment: 2$\times$ NVIDIA GeForce RTX 3090 (24\,GB), PyTorch 2.6.0+cu124, Python 3.12.3 (\texttt{run\_manifest.json}).

Legacy benchmark enforced Transformers $\geq$4.51.0 and 4-bit QLoRA loading~\cite{dettmers2023qlora}.
Training follows parameter-efficient adaptation rather than full-model fine-tuning~\cite{howard2018ulmfit,hu2022lora}.

\section{Legacy Benchmark (Phase A)}

Script \texttt{run\_skillspan\_benchmark.sh} trains and evaluates five models sequentially with seed 42, batch size 1, temperature 0.
Mistral legacy config: 3 epochs, LoRA rank 64, learning rate $1.2 \times 10^{-4}$, 4-bit quantization (\texttt{configs/mistral\_small\_3\_1\_v2.yaml}).

\section{Stage-2 Training (Phase B)}

Mistral LoRA: 80 epochs, LoRA rank 16, alpha 32, batch size 1, gradient accumulation 32, max sequence length 512, 2-GPU DDP (\texttt{skillllm\_v2\_training\_80ep\_2gpu.yaml}).
Training runtime: 324{,}218 seconds ($\approx$90 hours); train loss 0.143; dev eval loss 0.397 at epoch 80 (\texttt{training\_summary.json}).

Single-run evaluation; multi-seed variance was not recorded in artifacts.

% ============================================================
\chapter{Results}
\label{ch:results}

\section{Phase A: Multi-LLM Benchmark}

Table~\ref{tab:legacy} reports completed three-epoch results from \texttt{summary\_report.json}.
Gemma 4 26B did not complete (meta-tensor and CUDA OOM errors in \texttt{errors.jsonl}).

\begin{table}[htbp]
  \centering
  \caption{Legacy SkillSpan benchmark (3 epochs, 3{,}569 test examples).}
  \label{tab:legacy}
  \begin{tabular}{lrrr}
    \toprule
    Model & Exact F1 & Token F1 & Partial F1 \\
    \midrule
    Mistral Small 3.1 24B & 0.665 & 0.730 & 0.798 \\
    Qwen3 14B            & 0.649 & 0.714 & 0.789 \\
    Gemma 4 E4B           & 0.636 & 0.705 & 0.781 \\
    Llama 3.1 8B Instruct & 0.632 & 0.704 & 0.764 \\
    Gemma 4 26B A4B       & ---   & ---   & --- \\
    \bottomrule
  \end{tabular}
\end{table}

Legacy Mistral per-label F1: hard 0.730, soft 0.587.
Parse failure rate: 0.028\% (1/3{,}569) for each completed model.

\section{Phase B: Stage-2 Extraction}

Table~\ref{tab:stage2} reports verified test metrics from Stage-2 artifact JSON files.

\begin{table}[htbp]
  \centering
  \caption{Stage-2 SkillSpan test metrics.}
  \label{tab:stage2}
  \begin{tabular}{lrrrrr}
    \toprule
    System & Exact & Token & Partial & Hard & Soft \\
    \midrule
    BERT baseline       & 0.059 & 0.275 & 0.176 & 0.130 & 0.019 \\
    GLiNER (zero-shot)  & 0.076 & 0.114 & 0.135 & 0.055 & 0.097 \\
    Raw Mistral (80 ep.)& 0.670 & 0.788 & 0.814 & 0.726 & 0.604 \\
    Full CAG/graph wrap.& 0.671 & 0.788 & 0.816 & 0.729 & 0.604 \\
    \bottomrule
  \end{tabular}
\end{table}

Extended Mistral training improved exact span F1 from 0.665 to 0.670 absolute over Phase A.
The full wrapper adds taxonomy linking coverage 0.890 with gate decision \texttt{accept} and exact span F1 delta +0.001 versus the anchor.

\section{Downstream Metrics}

\texttt{matching\_eval.csv} and \texttt{course\_recommendation\_eval.csv} record NA for nDCG@10 and P@5.
Demo JSON outputs illustrate job match scores on synthetic inputs only.

\section{Smoke Ablation}

A 20-example smoke ablation (\texttt{extraction\_ablation.md}) reports much lower F1 (e.g., 0.057 for Mistral-only).
Those numbers reflect the smoke sample and are not used for primary conclusions in this draft.

% ============================================================
\chapter{Discussion and Error Analysis}
\label{ch:discuss}

\section{Model Selection}

Among four completed LLMs at equal three-epoch budget, Mistral leads on all primary span metrics despite higher latency than Llama 3.1 8B (2.03\,s vs.\ 0.91\,s per example in legacy logs).
Capacity and instruction-following appear to outweigh raw inference speed for boundary-sensitive extraction.

\section{Generative vs.\ Encoder Approaches}

Mistral exact span F1 exceeds BERT~\cite{devlin2019bert} and zero-shot GLiNER~\cite{zaratiana2023gliner} by an order of magnitude under shared Stage-2 evaluation.
JSON generative extraction with supervised LoRA aligns with SkillSpan's span-and-label interface and with recent LLM extraction findings~\cite{nguyen2024rethinking,herandi2024skillllm}.
The repository's earlier NNOSE baseline on resume sentences achieved lower span F1 under a different protocol~\cite{zhang2024nnose}, reinforcing that benchmark parity matters when comparing systems.

\section{Wrapper Impact}

Wrapper components yield incremental extraction gains (+0.001 exact span F1 for full CAG/graph) but add taxonomy linking at 0.890 coverage.
The primary Stage-2 engineering contribution is integrated, traceable skill intelligence with non-regression guards rather than a large F1 breakthrough.

Wrapper error analysis on 974 test documents reports two \texttt{bad\_boundary\_edit} regressions attributed to the taxonomy linker in GLiNER+taxonomy and full graph variants.
Partial span F1 consistently exceeds exact span F1 (0.814 vs.\ 0.670), indicating boundary near-misses as the dominant residual error class.
Soft-skill F1 trails hard-skill F1 by approximately 0.12 points across strong models, consistent with prior soft-skill matching and annotation studies~\cite{sayfullina2018softskill,calanca2019responsible}.

% ============================================================
\chapter{Limitations, Reproducibility, and Future Work}
\label{ch:limits}

\section{Limitations}

\begin{itemize}
  \item SkillSpan domains may not generalize to resumes or multilingual labor markets~\cite{vasquezrodriguez2026sixlanguages,kavas2025multilingual}.
  \item Gemma 4 26B was not evaluable on 2$\times$24\,GB GPUs.
  \item Phase A and Phase B differ in epoch count and LoRA configuration.
  \item Single-seed runs; small wrapper deltas should be interpreted cautiously.
  \item No labeled job/course evaluation for downstream modules.
  \item This LaTeX source has not passed SFSU preliminary format check.
\end{itemize}

\section{Reproducibility}

Legacy benchmark:
\begin{verbatim}
bash scripts/run_skillspan_benchmark.sh
\end{verbatim}
from \texttt{SkillLLM\_v2/}.

Stage-2:
\begin{verbatim}
bash scripts/run_skillllm_v2_smoke.sh
PYTHONPATH=src pytest tests/skillllm_v2/ -q
\end{verbatim}
from \texttt{SkillLLM\_v2\_stage2/}.

Master repo commit: \texttt{0627abf5170ccbe5fe01e5ab1238dc5ac340b880}.

\section{Future Work}

Construct labeled job/course evaluation sets aligned with ConFit-style resume--job matching benchmarks~\cite{yu2025confitv2,vaishampayan2025humanllm}; run multi-seed training; complete trained span verifiers~\cite{ma2023llminformationextractor}; migrate graph layer to production stores; evaluate bias in matching~\cite{iso2025bias}; integrate live job ingestion under privacy constraints.

% ============================================================
\chapter{Conclusion}
\label{ch:concl}

This thesis documented SkillLLM\_v2 from repository evidence: a SkillSpan benchmark selecting Mistral Small 3.1 24B among feasible open-weight models, followed by a context-aware graph-augmented Stage-2 system with verified extraction metrics, taxonomy linking, and guarded wrapper design.
Encoder baselines remained far behind generative Mistral extraction.
Downstream recommendation requires labeled evaluation before quantitative claims can be made.
Soft-skill extraction and span boundaries remain open problems for future work.

% ============================================================
\bibliographystyle{plain}
\bibliography{references}

\end{document}
